\section{Existing Literature on Household Credit and Buy Now, Pay Later}
\label{ch:literature}

Four literatures provide foundation for the research questions: agent-based models of household credit supply
the method and its validation standard, the empirical \ac{BNPL} record the phenomenon, the
behavioural evidence the decision rules, and the regulatory record both the institutional setting
and the interventions tested under the third research question.

\subsection{Agent-Based Models of Household Credit}
\label{sec:lit-abm}

Distress at the level of a population is more than the sum of distress at the level of a
household, which is the argument for modelling it from the bottom up. Cardaci
\cite{cardaci2018inequality} builds a hybrid agent-based stock-flow-consistent model in which
rising inequality, transmitted through expenditure cascades, pushes households into debt-financed
consumption until accumulating defaults produce an endogenous banking crisis, with severity
governed by credit-market conditions and the strength of consumption emulation. Their household debt
is collateralised, structurally closer to mortgages than to unsecured pay-in-four instalments,
a difference noted wherever their mechanisms are applied here. D'Orazio and Giulioni
\cite{dorazio2017micro} reach a similar conclusion: unsustainable debt
emerges from local borrowing rules rather than being written into the model. Theese findings justify the
methodological premise for this thesis, since the mechanism under investigation, a single household
holding obligations to several lenders none of whom can observe the others, exists only on the
individual balance sheet and cannot be expressed by a representative agent.

Madeira \cite{madeira2018chile}, at the Central Bank of Chile, models households that absorb
labour-income shocks, choose among loan contracts, and default when they can no longer finance
minimum consumption. That default condition is a cash-flow insolvency test rather than a leverage
threshold, and it is adopted here because the \ac{NCA} affordability assessment is itself a
residual-income test. Madeira also validates against observed Chilean default rates, demonstrating
that a household-level credit \ac{ABM} in a middle-income economy can be held to external
behavioural targets, which is the standard this thesis adopts.

Hamill et al. \cite{hamill2023creditcard} model the United Kingdom credit-card market and find
that promotional strategies chosen by profit-maximising lenders raise borrower indebtedness. Their
model also reproduces a non-monotonic relationship between income and debt-to-income ratio, with
middle-income households carrying the highest values. Because that profile comes from a model
rather than from data, it serves here as a cross-model plausibility check and not as an empirical
validation target.

Windrum, Fagiolo and Moneta \cite{windrum2007validation}
observe that \ac{ABM} parameters may grow unbounded when seeking to fit the model to data, and that a model
judged on the patterns it was fitted to has limited predictive power. The methodology adopted in this thesis,
following the pattern-oriented approach embedded in the \ac{ODD} protocol \cite{grimm2020odd,grimm2006odd}, is
to register empirical patterns in advance, fit as few parameters as possible, and judge the model
on patterns it was never fitted to.

\subsection{Empirical Evidence on Buy Now, Pay Later}
\label{sec:lit-bnpl}
deHaan et al. \cite{dehaan2024bnpl} provide evidence related to the banking
records introduced in Section~\ref{sec:background}, new \ac{BNPL} users incur higher
overdraft charges, credit-card interest and late fees than comparable non-users, with the
deterioration concentrated among the liquidity-constrained. A contrasting view, that \ac{BNPL}
substitutes for traditional credit is unsupported, Di Maggio et al. \cite{dimaggio2022bnpl}, show that \ac{BNPL} access raises total
spending by more than intertemporal substitution can explain, for constrained and unconstrained
users alike. These findings provide the empirical basis that this model aims to reproduce.

The \ac{CFPB} \cite{cfpb2025bnpl} documents a mechanism by which household exposure grows. They find that
two-thirds of observed borrowers held simultaneous loans at some point and a third borrowed from more than
one firm at once, across providers that cannot see one another's ledgers, so an interest-free
deferral on a single transaction can accumulate into an aggregate servicing burden that no
individual lender ever assessed. Hayashi and Routh
\cite{hayashi2025constraints} find \ac{BNPL} users markedly more likely to be financially
constrained than non-users on Federal Reserve survey measures, and Wang \cite{wang2026buynow}
reports use concentrated among lower-income and lower-credit-score consumers. These findings justify
the model's focus on the lower-income quintiles and on the consequences of concurrent-facility stacking.

On the demand side, Ackert et al. \cite{ackert2025bnpl} ran an incentivised experiment in which
participants were more likely to finance a purchase with \ac{BNPL} than with a credit card, and
most believed their peers would approve; they conclude that indebtedness is likely to rise when a
payment scheme carries that social sanction.
Granovetter \cite{granovetter1978threshold} shows that share-dependent adoption with
heterogeneous individual thresholds can produce discontinuous aggregate outcomes, which indicates that
models of \ac{BNPL} adoption should include a social component to capture the observed rapid uptake. These findings provide the rationale for
the inclusion of peer-effects and want-driven borrowing within the model.



\subsection{Behavioural Foundations for the Decision Rules}
\label{sec:lit-behavioural}

Experimentally elicited present bias is a strong predictor of credit-card borrowing
\cite{meier2010present}, and present-biased borrowers systematically fail to execute the debt
paydowns they plan \cite{Kuchler2021}. Together these justify a borrowing trigger that reacts to
an immediate shortfall rather than to a lifetime budget constraint, and repayment behaviour that
cannot be assumed to follow the contractual schedule.

Keys and Wang \cite{Keys2019} demonstrate that 29\% of United States credit-card
accounts surveyed regularly pay at or near the contractual minimum, and 9 to 20\% of accounts respond to
changes in the minimum-payment formula in ways liquidity constraints alone cannot explain, a lower
bound on anchoring to the stated minimum. Uniform scheduled amortisation is therefore rejected as
a universal repayment rule within the model. The model instead assigns each
household one of two fixed repayment archetypes, minimum-payer or scheduled-payer, with the
minimum-payer share set at the observed 0.29 and swept around it; a continuous present-bias
parameter was considered and rejected because no South African distribution exists to calibrate
one.

All of this evidence is from the United States. Whether parameters estimated on credit-card
borrowers in a high-income economy transfer to a population in which grant income dominates the
bottom quintiles is untested, and Section~\ref{sec:lim-transfer} treats that transfer as a
limitation in its own right.

\subsection{The Regulatory Setting, at Home and Abroad}
\label{sec:lit-regulatory}

The statutory architecture described in Section~\ref{sec:background} rests on the \ac{NCA}
\cite{sa_nca_2005} and the Regulation 23A residual-income test \cite{sa_ncr_affordability_2015},
and \ac{BNPL} providers position their product outside both \cite{nortonrose_bnpl_sa}. One
further provision matters for the model: since May 2016 the \ac{NCA} has capped the maximum rate
on unsecured transactions at the repurchase rate plus 21\% \cite{sa_ncr_ratecaps_2016}, which
sets the cost boundaries for the traditional credit products in the model.

Local measurement of the resulting exposure is thin. TransUnion's consumer survey
\cite{transunion_cps_sa_2025} establishes that \ac{BNPL} adoption amongst surveyed South African consumers is increasing,
noting that \ac{BNPL} facilities are highly sought after. However, self-reported holding in a marketing survey is weak evidence of exposure levels,
but it establishes that adoption is not marginal, and the same source confirms that \ac{BNPL} payment histories remain absent from
South African credit reports. That invisibility is structural rather than accidental, which is
why bureau visibility enters the model as a switchable policy lever rather than a constant.

Regulators elsewhere have already moved against the same gap, and their instruments define the
intervention levers used in this model. In the United Kingdom, the Woolard Review
\cite{woolard2021} initiated the process that produced the \ac{FCA}'s final rules
\cite{fca_ps26_1}, which bring deferred payment credit under proportionate mandatory
affordability checks and a fourteen-day right of withdrawal. In
Australia, \ac{ASIC}'s industry review \cite{asic2020rep672} found that one in five surveyed users
had missed a payment in the prior year under industry self-regulation. Affordability checks,
cooling-off periods and reporting obligations are therefore not hypothetical constructs but the
actual policy repertoire, and the third research question asks which of them works in the South
African market structure.

\subsection{Synthesis and Gap}
\label{sec:lit-gap}

Table~\ref{tab:gap}
summarises the findings of the previously described papers, their contributions and gaps.

\begin{table}[H]
\centering
\caption{What each literature establishes, and what it leaves open.}
\label{tab:gap}
\begin{tabular}{L{3.2cm}L{5.0cm}L{5.2cm}}
\toprule
\textbf{Literature} & \textbf{Establishes} & \textbf{Leaves open} \\
\midrule
Empirical \ac{BNPL} \newline \cite{dehaan2024bnpl, cfpb2025bnpl, dimaggio2022bnpl,
hayashi2025constraints, wang2026buynow} & That \ac{BNPL} adoption is followed by worse financial
outcomes, that spending rises rather than substitutes, and that concurrent multi-firm borrowing is
widespread. & All from United States data, where bureaus partly observe \ac{BNPL}. Outcomes are
measured; the mechanism by which invisible exposures aggregate is not. \\
\addlinespace
Behavioural \newline \cite{meier2010present, Keys2019, Kuchler2021, ackert2025bnpl,
granovetter1978threshold} & The individual drivers of debt: shortfall borrowing, payment
anchoring, and socially sanctioned adoption. & The drivers are measured one consumer at a time; an
explicit aggregation structure is needed to connect them to market-level outcomes, which is what
this model adds. \\
\addlinespace
Household credit \ac{ABM} \newline \cite{cardaci2018inequality, dorazio2017micro,
madeira2018chile, hamill2023creditcard, windrum2007validation} & That individual credit decisions
can be simulated to macro-level instability, and the out-of-sample standard such models must
meet. & All assume a single regulatory perimeter. None reviewed carries a lender class exempt
from both affordability checks and bureau visibility. \\
\addlinespace
Regulatory \newline \cite{sa_nca_2005, sa_ncr_affordability_2015, sa_ncr_ratecaps_2016,
nortonrose_bnpl_sa, transunion_cps_sa_2025} & A bifurcated South African market in which
registered lenders share data while \ac{BNPL} providers do not, and the intervention instruments
adopted abroad. & The household-level consequences: administrative data cannot record what is
never reported. \\
\bottomrule
\end{tabular}
\end{table}